\section{Restriction for the fixed-conductor characters}
\label{iv-section}

We next prove the arithmetic assertion called case $IV$ in the
submitted construction.\cite{alpo-uniform-source} The main point is
that its character has conductor dividing the original fixed modulus.
Consequently it is trivial over one of the fixed odd-prime fields
already controlled by the common cohomological interval of
Section~\ref{sec:selection-assembly}. We keep the resulting bound
additive in that interval's bound $B$.

\subsection{The character loses its varying conductor}

Fix a primitive Dirichlet character $\chi$ of conductor $M$ and
order $o$, and retain the fields
\[
 E_h=K_\chi R_h(\sqrt{-q},\sqrt{-d}),\qquad
 N_h=\operatorname{lcm}(2^{h+2},M)qd
                                                               \tag{IV1}
\]
from Section~\ref{sec:selection-assembly}. Thus $q,d$ are distinct
odd primes prime to $M$, and $R_h$ is the real subfield of the
$2^{h+2}$-th cyclotomic field. A character $\theta$ of
$\operatorname{Gal}(E_h/\mathbb Q)$ will also denote its
Dirichlet character modulo $N_h$.

For a prime $r$, write $\theta_r$ for the component on
$(\mathbb Z/r^{v_r(N_h)}\mathbb Z)^\times$ under the Chinese
remainder decomposition; it is understood to be the trivial
character when $r\nmid N_h$. Independently, for an odd prime
$p$, write
\[
                    \theta=\kappa\vartheta,
\]
where $\kappa$ has $p$-power order and $\vartheta$ has
order prime to $p$. This factorization is unique in the finite
cyclic group generated by $\theta$. Define
\[
                              \eta=\vartheta/\vartheta_p.
                                                               \tag{IV2}
\]
The denominator in (IV2) is the local component at the rational
prime $p$; it is distinct from the factor of $p$-power order.

\begin{lemma}\label{iv-fixed-conductor}
Suppose $p\ne q,d$ and $\eta$ has odd order. Then
\[
                              \operatorname{cond}\theta\mid M.
                                                               \tag{IV3}
\]
Indeed the local components of $\theta$ at $2,q,d$ are trivial.
\end{lemma}
\begin{proof}
The group of units modulo a power of $2$ is a $2$-group, so
$\theta_2$ has $2$-power order. Its $p$-power order factor is
trivial because $p$ is odd. Removing the local component at $p$
does not affect it. Hence $\eta_2=\theta_2$. Its order divides
both a power of $2$ and the odd order of $\eta$, and is therefore
one.

The inertia groups at $q$ and $d$ in $E_h/\mathbb Q$ have order
at most two. This follows directly from (IV1): the other fields
are unramified there, and the respective added quadratic field
has degree two. Equivalently, the two local Dirichlet components
are trivial or quadratic. Their $p$-power order parts are again
trivial. Since $p\ne q,d$, they are unchanged in (IV2), so their
orders are odd and divide two. They too are trivial.

At every remaining odd prime, the exponent in the modulus $N_h$
is precisely the exponent in $M$. Triviality at $2,q,d$ now
shows that the primitive conductor divides $M$.
\end{proof}

For the selected fixed set of odd primes $P$, the inequalities
$q,d>\max P$ ensure $p\ne q,d$ for every $p\in P$.
Thus the hypothesis of Lemma~\ref{iv-fixed-conductor} is exactly
the odd-order alternative required in case $IV$. No bound for
$q,d$ enters the conclusion.

\subsection{Restriction and the loss in degree one}

For a number field $F$ and $e\ge1$, let
\[
 \mathcal S_p(F,e)=\ker\left[
 H^1(G_F,\mathbb Q_p/\mathbb Z_p(-e))\longrightarrow
 \prod_{v\mid p}H^1(D_v,\mathbb Q_p/\mathbb Z_p(-e))
 \oplus\prod_{v\nmid p}H^1(I_v,\mathbb Q_p/\mathbb Z_p(-e))
 \right].                                                     \tag{IV4}
\]
All primes in (IV4) are finite. Since $p$ is odd, the corresponding
cohomology at a real place vanishes by averaging over its group
of order two. Put
$\lambda_p(F,e)=\operatorname{length}_{\mathbb Z_p}\mathcal S_p(F,e)$,
with value $+\infty$ when this length is infinite.

Let $K/\mathbb Q_p$ be finite, with ring of integers $\mathcal O$
and ramification index $b$. Suppose that the values of
$\theta$ belong to $\mathcal O^\times$. Put
\[
 A(e,\theta)=(K/\mathcal O)(\theta^{-1}\varepsilon_p^{-e}).
                                                               \tag{IV5}
\]
Let $\mathcal T_p(e,\theta)$ be the strict Selmer group defined
as in (SZ2), with this coefficient module, and write
\[
 s_p(e,\theta)=b^{-1}
       \operatorname{length}_{\mathcal O}\mathcal T_p(e,\theta)
                                                               \tag{IV6}
\]
when the length is finite. Additional zero local conditions at
primes dividing $N_h$ only decrease this length.

\begin{lemma}\label{iv-restriction}
Suppose $\operatorname{cond}\theta\mid M$. For an integer
$j\ge0$, set
\[
 F=\mathbb Q(\mu_{p^{j+1}M}),\quad
 a=j+1+v_p(M),\quad n=[F:\mathbb Q].
                                                               \tag{IV7}
\]
If $\mathcal S_p(F,e)$ is finite, then $\mathcal T_p(e,\theta)$
is finite and
\[
             s_p(e,\theta)\le\lambda_p(F,e)+n(a+v_p(e)).
                                                               \tag{IV8}
\]
In particular this statement permits $p\mid n$ and arbitrary
ramification in $K/\mathbb Q_p$.
\end{lemma}
\begin{proof}
The character $\theta$ is trivial on $G_F$. Let
$\Delta=\operatorname{Gal}(F/\mathbb Q)$ and
$A=A(e,\theta)$. Inflation--restriction gives an exact sequence
\[
 0\longrightarrow H^1(\Delta,A^{G_F})
 \longrightarrow H^1(G_{\mathbb Q},A)
 \longrightarrow H^1(G_F,A).
                                                               \tag{IV9}
\]
No division by $|\Delta|$ is used. A class in
$\mathcal T_p(e,\theta)$ restricts into the strict Selmer group
over $F$: a decomposition or inertia group over $F$ is a
subgroup of the corresponding conjugate group over $\mathbb Q$.
The local restriction of a coboundary is still a coboundary.

On $G_F$ the module $A$ is $(K/\mathcal O)(-e)$. Its strict
Selmer group is
\[
                       \mathcal S_p(F,e)\otimes_{\mathbb Z_p}
                                      \mathcal O.              \tag{IV10}
\]
To see this directly, choose a $\mathbb Z_p$-basis of the finite
free module $\mathcal O$. Since
$K/\mathcal O=(\mathbb Q_p/\mathbb Z_p)\otimes\mathcal O$,
the cochain complex, all local cochain complexes, and their
restriction kernels become the direct sums indexed by that
basis. The resulting identification is the canonical
$\mathcal O$-linear map in (IV10). A cyclic summand
$\mathbb Z_p/p^r\mathbb Z_p$ contributes
$\mathcal O/p^r\mathcal O$, of $\mathcal O$-length $br$.
The target of restriction therefore has length
$b\lambda_p(F,e)$.

It remains to bound the left group in (IV9). The cyclotomic
image of $G_F$ is exactly $1+p^a\mathbb Z_p$. Indeed the
Chinese remainder theorem gives
\[
 F\cap\mathbb Q(\mu_{p^\infty})=\mathbb Q(\mu_{p^a}).
\]
For odd $p$ and $a\ge1$,
\[
                   v_p((1+p^a)^e-1)=a+v_p(e).                 \tag{IV11}
\]
For completeness, when $p\nmid e$ the first nonconstant term
of the binomial expansion has valuation $a$ and every other
term has larger valuation. If $v_p(z)\ge1$, the term $pz$
in $(1+z)^p-1$ has valuation $1+v_p(z)$, strictly below
all the other terms, because $p$ is odd. Iterating this
observation proves (IV11) for every $e\ge1$.
The same formula holds with $e$ replaced by $-e$, since
$(1+p^a)^e$ is a unit. Every element of
$1+p^a\mathbb Z_p$ has its $e$-th power congruent to one
modulo $p^{a+v_p(e)}$, while $1+p^a$ attains this valuation.
It follows that
\[
 A^{G_F}=A[p^{a+v_p(e)}],\qquad
 \operatorname{length}_{\mathcal O}A^{G_F}=b(a+v_p(e)).
                                                               \tag{IV12}
\]
Every $1$-cocycle of the finite group $\Delta$ is a function
$\Delta\to A^{G_F}$. Thus its cocycle module is a submodule
of $(A^{G_F})^n$, and taking the quotient by coboundaries
cannot increase length. Consequently
\[
 \operatorname{length}_{\mathcal O}H^1(\Delta,A^{G_F})
                                  \le nb(a+v_p(e)).
\]
Combining this with (IV9)--(IV10) proves finiteness and (IV8).
\end{proof}

\subsection{Summation with fixed constants}

\begin{theorem}\label{iv-summed}
Fix $p,M,j$, and let $F,a,n$ be as in (IV7). If $t\ge2$
and
\[
                 \sum_{1\le e<t}\lambda_p(F,e)\le B,
                                                               \tag{IV13}
\]
then for every character $\theta$ of conductor dividing $M$
all the corresponding strict Selmer groups are finite, and
\[
 \sum_{1\le e<t}s_p(e,\theta)
 \le B+n\left(a+\frac1{p-1}\right)(t-1).
                                                               \tag{IV14}
\]
There is an integer $C$, depending only on $p,M,j$, such that
for every $N\ge1$ the right side is at most $B+(Nt)^C$.
This bound remains true after adding the level-$N$ zero
local conditions.
\end{theorem}
\begin{proof}
Sum (IV8). Counting the multiples of each power of $p$ gives
\[
 \sum_{e=1}^{t-1}v_p(e)
   =\sum_{r\ge1}\left\lfloor\frac{t-1}{p^r}\right\rfloor
   \le\frac{t-1}{p-1}.
\]
This proves (IV14). Choose $C\ge1$ so that
$2^{C-1}\ge n(a+1)$. Since $t\ge2$ and $N\ge1$,
$n(a+1)t\le t^C\le(Nt)^C$.
The last assertion follows from inclusion of Selmer groups.
\end{proof}

\begin{corollary}\label{iv-exact-case}
Fix the original character $\chi$, an odd prime $p$ in the
fixed set $P$, and the fixed cyclotomic layer index $j_p$.
Use the common interval with base field
$\mathbb Q(\mu_{pM})$ at $p$. For every growing field (IV1),
every character $\theta$ in its case $IV$, and every $t$
controlled by that interval, the normalized sum of its
level-$N_h$ Selmer lengths is at most
\[
                              B+(N_ht)^C,
                                                               \tag{IV15}
\]
where $C$ depends only on $\chi,p,j_p$.
In particular $C$ is independent of $h,q,d,t$ and of the
chosen coefficient extension.
\end{corollary}
\begin{proof}
The $j_p$-th layer of the cyclotomic $\mathbb Z_p$-extension
of $\mathbb Q(\mu_{pM})$ is
$\mathbb Q(\mu_{p^{j_p+1}M})$: its degree over the base is
$p^{j_p}$. These are precisely the fixed fields entering
(IV13). The prime selection makes $q,d>\max P$.
Lemma~\ref{iv-fixed-conductor} gives
$\operatorname{cond}\theta\mid M$, and
Theorem~\ref{iv-summed} gives (IV15).
\end{proof}

There are at most $\varphi(M)$ distinct characters of conductor
dividing $M$. Thus, if (IV13) holds with a common $B$ for
$p\in P$, the entire collection just treated contributes at
most
\[
 \varphi(M)\left(\sum_{p\in P}\log p\right)B
  +\varphi(M)\sum_{p\in P}\log p\,
       n_p\left(a_p+\frac1{p-1}\right)(t-1)                 \tag{IV16}
\]
to the sums of character lengths weighted by $\log p$.
Both coefficients in (IV16) are fixed before selecting $q,d$
and the growing $2$-power conductor. This establishes the
required dependence on $B$ for these characters; it does not
replace the separate estimates for the other characters.


